\centerline{\bf \Huge Всеросс}

\begin{flushright}
        \scriptsize{Каждая задача абсолютно решаема}
\end{flushright}

\problem{1}
Каждая из функций $f(x)$ и $g(x)$ определена на всей числовой прямой и не является строго монотонной. Может ли быть, что и их сумма, и их разность строго монотонны на всей числовой прямой?
%УТГ 2021

\problem{2}
Возрастающая последовательность натуральных чисел $a_1<a_2<\ldots$ такова, что при каждом целом $n>100$ число $a_n$ равно наименьшему натуральному числу, большему чем $a_{n-1} u$ не делящемуся ни на одно из чисел $a_1, a_2, \ldots, a_{n-1}$. Докажите, что в такой последовательности лишь конечное количество составных чисел.
%УТГ 2021

\problem{3}
На прямоугольном столе лежат равные картонные квадраты $n$ различных цветов со сторонами, параллельными сторонам стола. Если рассмотреть любые $n$ квадратов различных цветов, то какие-нибудь два из них можно прибить к столу одним гвоздем. Докажите, что все квадраты некоторого цвета можно прибить к столу $2 n-2$ гвоздями.
%Финал 2000

\problem{4}
Неравнобедренный треугольник $A B C$ вписан в окружность с центром $O$ и описан около окружности с центром $I$. Точка $B^{\prime}$, симметричная точке $B$ относительно прямой $O I$, лежит внутри угла $A B I$. Докажите, что касательные к окружности, описанной около треугольника $B B^{\prime} I$, проведенные в точках $B^{\prime}$ и $I$, пересекаются на прямой $A C$.
%Финал 2016-2017